\documentclass[11pt]{article}
\usepackage{lmodern}
\usepackage[T1]{fontenc}
\usepackage[letterpaper,margin=0.8in]{geometry}
\usepackage{xcolor}
\usepackage{booktabs}
\usepackage{tabularx}
\usepackage{enumitem}
\usepackage{array}
\usepackage[hidelinks]{hyperref}

\setlist{itemsep=1pt,topsep=3pt}
\newcommand{\h}[1]{\texttt{#1}}
\newcolumntype{L}{>{\raggedright\arraybackslash}X}

\title{LArPix ASICs on PACMAN\\[2pt]
  \large Bench setup, configuration, replies and status}
\author{Maintained by the PACMAN agents (augie)}
\date{2026-11-01}

\begin{document}
\maketitle
\sloppy

\section{Purpose and scope}

This note records the LArPix ASIC work on PACMAN with the zipline firmware and software (\h{zipline.tex}): which
ASICs are on the bench, how they are configured, what they reply, and what is open.  Register numbers are those of
the LArPix v3 register map in \h{larpix-lib} (\h{config/asics/larpix\_v3.yaml}), which agrees with the LBNL datasheet
(rev 2.4, not public).  The datasheet has known errors: its bit layouts and defaults are not trusted until they are
read back from a chip.

\section{Bench setup}

\begin{itemize}
\item One LArPix v3 on a single-tile board, on tile 1.  On the 1v4 it is reached through the J1302 cable, with all
  four PISO/POSI pairs routed to the tile-1 UARTs 0--3.
\item Chip ID 0x0B (11).  The hello request is a configuration read of register 122 (the chip ID); the expected reply
  is \h{0x42254139 1C2DE82F} (``read, chip 0x0B, register 122, value 0x0B'').
\item Tile power by \h{pacman\_menu} (VDDA about 2.32\,V, VDDD about 1.20--1.30\,V, a few mA each).  The 1v4 needs I2C
  at 100\,kHz with this cable (now the default).
\item Unplug any signal source from the LEMO inputs before connecting ASICs.
\item After a full reset, wait (0.2\,s is enough) before configuring: the reset returns at once, and a configuration
  sent during it is lost.  A pollable reset state register is set aside for later.
\end{itemize}

\section{Configuration root}

The configuration root (\h{pacman\_menu}, \h{asic\_config\_root} in \h{asic\_menu.c}) brings a chip up as the root of
a chain:

\begin{center}
\begin{tabularx}{\linewidth}{@{}llL@{}}
\toprule
Register & Value & Meaning \\
\midrule
125 & 0x1 & downstream PISO enables: PISO 0 \\
124 & 0x2 & upstream PISO enables: PISO 1 \\
239, 240 & 0xFF & PISO drive: every TX slice enabled \\
241, 242 & 0x00 & PISO drive current (\h{i\_tx\_diff}, active low): full \\
\bottomrule
\end{tabularx}
\end{center}

\noindent The drive settings were added on 2026-10-31 (pacman-software 6d2ad8e) after the 1v4 results below.

\section{Results}

\begin{center}
\begin{tabularx}{\linewidth}{@{}llL@{}}
\toprule
Date & Board & Result \\
\midrule
2026-10-30 & 1v5 & Hello on UART 0 at the power-on drive: the expected reply, exactly.  POSI request and PISO reply also
  decoded bit-exact from oscilloscope captures at the tile connector. \\
2026-10-30 & 1v4 & At the power-on drive, the same wrong word on all four lanes (\h{0xFE000000 70387FC0}): the reply framed
  7 bits late, with bit errors.  All TX slices (239, 240 = 0xFF): exact on UARTs 0 and 1.  Plus full current (241,
  242 = 0x00): exact on all four. \\
2026-10-31 & 1v4 & With 124 = 0x2 and full drive: exact on UART 0, and the built-in logic analyzer shows the request on
  POSI 0 and the exact reply on PISO 0, 8.2\,\textmu s later.  PISO 2 and 3 (enabled neither up- nor downstream)
  stream junk into the receivers at full drive. \\
2026-10-31 & 1v4 & Latest build (firmware 928fcf0, software 3c9b2d0), configuration root as built: the reply exact,
  twice. \\
2026-11-01 & 1v4 & Software 671a3dd, the same sequence as a piped script: no reply (UART 0 streams junk).  Cause: the
  configuration root sent right after the full reset is lost; the full reset returns at once, and the new menu input
  no longer sleeps between commands.  With a 0.2\,s pause after the full reset: the reply exact, twice. \\
2026-11-01 & 1v4 & Software 290979d, a 16-BD DMA ring set up from the menu, with the 0.2\,s pause: the reply exact,
  twice. \\
2026-11-01 & 1v4 & Drive scan with the ASIC loopback test (800 readbacks of random register values per point): at
  the power-on current (241 = 0xCC) no slice setting (239 = 0x11 to 0xFF) passes; with all slices on, the current
  passes from 241 = 0x88 (0 errors) to 0x00, 0x99 is marginal (22 errors), 0xAA and weaker fail; one timeout at
  0x22. \\
\bottomrule
\end{tabularx}
\end{center}

\section{Open}

\begin{itemize}
\item Systematic adjustment of the PISO drive (registers 239--242, and \h{v\_cm\_lvds\_tx} 254--255) and of the PISO
  enables (124, 125): why the 1v4 needs more drive than the 1v5, why UARTs 2 and 3 need more than 0 and 1, and
  whether disabled PISOs at full drive are undriven or biased without data.
\item The real bit layouts and defaults: a read-all of the registers on the 1v5.
\item ASIC stress tests on the 1v5 (later).
\item LArPix v2b on the hw1v4: v2 clocks its UART at MCLK/2 (v3 at MCLK).
\end{itemize}

\end{document}
